\documentclass[11pt]{article}
\usepackage{amsmath,amssymb,amsthm,mathtools,hyperref}

\begin{document}
\section*{Single-cell state dynamics}
\label{prob:048}
\label{rank:61}
\par\smallskip

\noindent\textbf{Domain.} Single-cell dynamics from snapshots. \quad \textbf{Retained status/source context.} Open constrained dynamical inference programme with inherent non-identifiability. Marginal snapshots alone do not determine a stochastic generator; optimal-transport methods add explicit assumptions and provide useful restricted reconstructions.
\paragraph{Definitions and assumptions.} Use a fixed cell-state representation in $\mathbb{R}^d$. Let $N_t(C)$ count living cells whose states lie in a measurable set $C$, and define the unnormalized expected cell-count density by $\mathbb{E}N_t(C)=\int_C p_t(x)\,dx$, measured with destructive state samples and absolute cell counts at specified times. The state $X_t$ in a transition probability refers to the cell or lineage sampling law declared below; $p_t$ is not its probability density by definition. A candidate diffusion has drift $b_\theta(x,t)$ and known or parameterized spatially constant symmetric positive semidefinite diffusion tensor $D_\theta$, with net growth/death rate $r_\theta(x,t)$. Measurement dropout and cell-state embedding are calibrated; birth/death rates or lineage information are measured when available. The dynamics satisfy $\partial_t p=-\nabla\cdot(bp)+\nabla\cdot(D\nabla p)+rp$ with no probability normalization; $\int p_t$ is expected population size. Division/lineage weighting and treatment of death are declared in transition validation. The declared offspring-state and death mechanism must be compatible with this population equation, including its local source term $rp$; arbitrary daughter-state redistribution is not represented by that term. The displayed $\mathcal L_\theta$ is the within-cell diffusion generator, not by itself a specification of branching, killing or lineage selection.
\paragraph{Problem statement.} Determine which transition predictions are identifiable and estimate them within the allowed generator class:
\[
\mathcal L_\theta a=b_\theta\cdot\nabla a+\operatorname{tr}(D_\theta\nabla^2a),
\quad |\widehat P(X_T\in A\mid X_0\in B)-P^{\rm val}(A\mid B)|\leq\varepsilon.
\]
Use independent lineage or live-cell validation for declared sets $A,B$. Here $P^{\rm val}(A\mid B)$ denotes the population transition probability under the declared validation sampling law, with the same lineage selection and conditioning convention as the prediction. A concrete instance declares the measurable set pairs, horizon $T$, positive tolerance $\varepsilon$ and whether coverage is pointwise or simultaneous over its validation domain. Each conditioning event has positive probability under the applicable law, or an explicit undefined/zero convention is declared; extinction and survival conditioning must be treated consistently. Finite validation estimates are distinguished from $P^{\rm val}$ and assessed using the binding uncertainty convention below, with dependence from shared ancestry or experimental replicates included in their sampling law. The practical goal is recovering reliable state transitions rather than merely matching density snapshots.
\paragraph{Additional binding assumptions and resolution criterion.} Distinguish a chronologically followed cell, a sampled surviving lineage and a size-biased descendant at time $T$; their transition laws generally differ under state-dependent proliferation. Specify which one is validated and how division affects state. The initial-cell sampling law, offspring-state and death rules, and lineage selection and conditioning rule must define that transition law jointly with the population model; a net rate $r_\theta$ alone does not specify these rules. The choice among these transition laws remains unspecified here. For fixed growth and diffusion, snapshot densities alone constrain a flux divergence, not a unique drift; if growth is unknown, transport and growth are constrained jointly. Lineage data or explicit restrictions are necessary to remove the remaining equivalences.

\paragraph{Alternative formulations.} (i) Assume deterministic gradient flow, known growth and sufficiently dense time samples, then infer a potential within a finite basis. (ii) Allow non-gradient drift and use lineage barcodes, velocity measurements or interventions to constrain probability currents. Adding a field $v$ with $\nabla\cdot(vp)=0$ can leave all marginals unchanged while changing trajectories, provided the perturbed drift remains in the candidate class and preserves the admitted solution and boundary or decay conditions, so unrestricted variants are not equivalent.
\paragraph{Obstructions and partial progress.} Waddington-OT supplies a concrete transport-based reconstruction with growth assumptions. Destructive samples do not pair ancestors and descendants; many couplings share the same marginals. Hidden cell-cycle states and incorrect growth rates bias trajectories. A valid answer may be bounds over compatible generators plus an experiment distinguishing them, rather than a unique inferred dynamic law.
\paragraph{Selective external check in this revision.} SC03 (13 October 2026; access level and exact support are recorded in the source-check appendix). These selected checks are not an exhaustive present-day status search.

\paragraph{Sources (supplied provenance; not all rechecked).}
\begin{itemize}
\item Fischer et al., Adapting systems biology to address the complexity of human disease in the single-cell era (2025; review metadata and abstract). \url{https://www.broadinstitute.org/publications/broad1361781}
\item Schiebinger et al., Optimal-Transport Analysis of Single-Cell Gene Expression Identifies Developmental Trajectories in Reprogramming (2019; primary assumption-dependent method). \url{https://pubmed.ncbi.nlm.nih.gov/30712874/}
\end{itemize}


\section*{Binding catalogue conventions}
This edition revises both supplied catalogues and preserves original problem numbers 1--330. It adds one globally unique label \verb|\label{rank:x}| per problem, where smaller $x$ means greater provisional difficulty. A rank label is an anchor at that problem's section: \verb|\ref{rank:x}| returns the section number, not the literal difficulty number. The visible \texttt{rank:x} text and the index give the difficulty number itself.

The entries are \emph{parameterized mathematical research schemes}, including empirical prediction, causal identification, historical reconstruction, theory and engineering. They are not 330 certified unsolved mathematical conjectures. A rigorous instance fixes the objects and restrictions below; changing those objects changes the question. Established identities, finite calculations, non-identification constructions and solved special cases are not being posed anew as open. Some residual tasks concern informative bounds, biological realization or external validation rather than a new theorem.

Every main definition and problem statement received a mathematical editorial review. Each entry now has an additional binding-conditions paragraph; direct corrections are separately recorded in the accompanying audit. The alternative formulations inherit the applicable conventions, but are different targets unless equivalence is proved. Where an old broad phrase conflicts with an explicit restriction below it, the explicit restriction controls.

\textbf{Literature boundary.} Thirteen selective primary-literature checks cover eighteen entries in this revision; their access levels and scope are recorded explicitly. The supplied bibliography and its earlier source-access descriptions are retained as provenance, not represented as newly verified. Neither this review nor a source's description proves the absence of later solutions. The date of this edition is 13 October 2026; there is no claim of exhaustive current-status certification for all 330 entries.

\subsection*{Binding conventions for a concrete instance}
\paragraph{Objects, quantifiers and non-vacuity.}
An instance declares measurable state, input, output and parameter spaces; a fixed candidate family; admissible initial laws, interventions and environments; a finite horizon or specified asymptotic regime; observation and sampling laws; and performance metrics, positive normalization scales and positive tolerances. These are chosen before fitting or testing, not adjusted after seeing failures. Shared parameters must genuinely be shared across the declared validation domain. Finite dimension alone does not restrict arbitrary real-number encodings: specify a finite library, bounded description length, or regularity/complexity bounds. Exclude constant or no-signal solutions through a fixed informative-output or baseline-improvement requirement where appropriate. A universally quantified result cannot be established by choosing one trivial favorable instance.

\paragraph{Probability laws and data.}
Every expectation or probability refers to the specified stochastic law, including initial uncertainty, process noise, observation noise and sampling as appropriate. A deterministic equation is not a probability model without a declared random input or initial law. ``Observed'' targets in population-level formulas mean population observables $g_j^*$, not noisy estimates asserted to be exact. Let $\widetilde g_j$ be independent validation estimates and suppose
\[
\mathbb{P}\{\lvert\widetilde g_j-g_j^*\rvert\le r_j\text{ for every }j\}\ge1-\alpha.
\]
Then an uncertainty-aware sufficient certification rule for positive target tolerances $\varepsilon_j$ is
\[
\max_j\frac{|\widehat g_j-\widetilde g_j|+r_j}{\varepsilon_j}\le1.
\]
On that simultaneous confidence event it implies every population error is at most its tolerance, by the triangle inequality. Vector/norm versions use corresponding valid simultaneous bounds. Training uncertainty must also be included when the claim averages over repeated training data. Evaluation on finitely many conditions does not establish a continuous-domain supremum without covering and regularity bounds.

\paragraph{Dynamics and well-posedness.}
For ODEs and SDEs give coefficient domains, initial laws and existence/uniqueness conditions in the intended solution sense, or keep the entire admitted solution set. For stochastic jump models specify nonnegative off-diagonal rates, diagonal normalization, rate orientation and non-explosion. For PDEs specify spatial domains, boundary/initial conditions, conservation or source terms, and a weak or classical solution class. Moving-domain concentration equations include dilution/transport terms. A stationary law is used only for a time-homogeneous augmented model under existence assumptions; a stationary approximation to a driven system must be justified. Absorbing extinction states are retained. Survival-conditioned, quasi-stationary and unconditional laws are distinct and must not be substituted for one another.

\paragraph{Identification and causal meaning.}
Two parameter values are observationally equivalent when they induce the same full distribution under all allowed observation/intervention designs, not merely the same likelihood value for one sample. A target is identified only if constant on the equivalence class. Gauge symmetries, label switching and reference conventions are quotiented or fixed. Finite-data compatible sets use a declared confidence, likelihood or discrepancy rule; an empty set signals model rejection. Bayesian credible sets are not automatically confidence sets. A prior or regularizer does not create data identification. Bounds intended to resolve a question must be informative at a predeclared width or decision threshold.

A causal contrast uses well-defined intervention versions and a stated causal model. Random assignment, or explicit consistency, overlap and exchangeability assumptions, is required for causal interpretation from observations. Interference and time-varying confounding are represented when present. Fixing a mediator by intervention differs from conditioning on its observed value. Marginal potential-outcome laws do not determine individual cross-world effects or an arbitrary joint coupling. Ethically admissible designs, rather than unconstrained interventions, define the experimental menu.

\paragraph{Metrics, units, optimization and limits.}
Specify the sampling or reference measure for densities and integrals, the log base, and dimensionless arguments of logarithms. Total variation is $\sup_A|P(A)-Q(A)|$; Wasserstein distances require the declared ground metric and appropriate moment conditions; KL divergence uses its stated absolute-continuity convention. All normed readouts are scaled consistently, and error probabilities use $0<\alpha<1$. A ratio has a positive denominator or an explicit undefined/zero convention. Thermodynamic entropy production is nonnegative in expectation under the relevant hypotheses; individual stochastic realizations need not be. An infimum is not an attained minimum without compactness/coercivity and continuity or another existence argument. Fix the order of large-population, long-time and vanishing-mutation limits. Stability claims declare perturbation space and symmetries.

\subsection*{How the provisional difficulty order was assigned}
Four editorial components are scored from 1 (lower burden) to 5 (higher burden): (a) formal or inferential difficulty, (b) access to evidence and identification burden, (c) mechanistic/computational integration, and (d) independent validation and transfer. Their unweighted total orders the entries. For selected entries in the highest bands, an explicit editorial priority breaks equal totals; all other ties are broken by original problem number. Equal totals are \emph{not} evidence of reliably different true difficulties. Scores are judgments, not measurements, completion-time estimates or consensus rankings.

The object ranked is an informative, nondegenerate resolution of the \emph{scoped} scientific target, with the stated budget and validation domain. A sharp useful impossibility result can qualify; an already known identity, vacuous bound, arbitrary favorable model or lookup table cannot. Rank 330 is merely the least difficult relative assessment in this catalogue, not a claim of an easy biological solution. Broad historical titles do not automatically outrank precisely formulated mechanistic tasks.


\section*{Selective source-check ledger}
\subsection*{SC03}\label{check:SC03}

Meng and Wang, Decoding gene regulatory networks from single-cell RNA velocity (arXiv preprint, 10 August 2026).

\textbf{Applied here to problems:} 46, 48, 72.

\textbf{Access level:} Primary preprint abstract and metadata; not an independent proof audit.

\textbf{Support and limit:} The preprint studies passive-data non-identifiability and restoration by perturbations under its model. It is not a universal identification guarantee for arbitrary latent biological networks.

\url{https://arxiv.org/abs/2608.09722}


\end{document}
